\documentclass[11pt]{article}
\usepackage[a4paper,margin=1in]{geometry}
\usepackage{amsmath,amssymb,mathtools,bm}
\usepackage{enumitem,booktabs,array}
\usepackage{tcolorbox}
\usepackage{hyperref}
\usepackage[protrusion=true,expansion=false]{microtype}
\hypersetup{colorlinks=true,linkcolor=blue,urlcolor=blue,citecolor=blue}
\setlist[itemize]{topsep=2pt,itemsep=2pt}
\setlist[enumerate]{topsep=2pt,itemsep=2pt}
\newtcolorbox{keybox}{colback=blue!3!white,colframe=blue!40!black,boxrule=0.4pt,arc=1mm}
\newtcolorbox{alertbox}{colback=red!3!white,colframe=red!40!black,boxrule=0.4pt,arc=1mm}
\newtcolorbox{answerbox}{colback=green!3!white,colframe=green!40!black,boxrule=0.4pt,arc=1mm}
\newtcolorbox{drillbox}{colback=yellow!5!white,colframe=orange!60!black,boxrule=0.4pt,arc=1mm}
\newcommand{\EV}{\mathbb{E}}
\newcommand{\Var}{\mathrm{Var}}
\newcommand{\Cov}{\mathrm{Cov}}
\newcommand{\blank}[1]{\underline{\hspace{#1}}}

\title{E06-02: Ahern \& Harford (2014, JF) \\
\large ``The Importance of Industry Links in Merger Waves'' --- Active Recall Workbook}
\author{Empirical Corporate Finance II --- Exam Prep}
\date{\today}
\begin{document}\maketitle

\begin{keybox}
\textbf{Paper in one sentence.} AH use BEA input-output tables to map supplier-customer links between industries and show that merger waves are not isolated industry events but \emph{propagate} through production networks: a wave in one industry predicts future waves in connected (upstream/downstream) industries, supporting a network-based theory where real-economy shocks cascade through the supply chain.

\textbf{Placement.} Key paper on M\&A-wave contagion; connects to Harford (2005) on wave origins. Complements network-based macro studies (Acemoglu et al.\ 2012).
\end{keybox}

\section*{Round 1: Complete Walk-Through}

\subsection*{1.1 Research question}
Harford (2005) establishes industry-level waves driven by shocks and liquidity. But industries are interconnected: do waves spread across supply chains? If yes, what explains the propagation?

\subsection*{1.2 Data}
\begin{itemize}
\item SDC mergers 1986--2010.
\item BEA benchmark input-output tables (multiple vintages): dollar flows from industry $i$ to industry $j$.
\item Compute industry-to-industry link strength: \% of $i$'s output going to $j$ (customer link) and \% of $i$'s input coming from $j$ (supplier link).
\end{itemize}

\subsection*{1.3 Empirical design}
\begin{itemize}
\item Industry-year panel of merger intensity (number of deals, aggregate value).
\item Lagged merger activity in customer/supplier industries as predictor.
\item Control for own-industry shocks, liquidity, and time effects.
\item Instrumental-variable design using technology diffusion patterns.
\end{itemize}

\subsection*{1.4 Headline results}
\begin{itemize}
\item A one-standard-deviation increase in lagged merger activity in customer/supplier industries predicts $\sim$20--30\% higher merger activity in a focal industry.
\item Upstream waves spread downstream more strongly than vice-versa.
\item Network spillovers operate over 1--3 year horizons; effect decays beyond that.
\item Consistent with consolidation propagating through the supply chain (e.g., telecom $\to$ equipment $\to$ component makers).
\end{itemize}

\subsection*{1.5 Mechanism}
\begin{itemize}
\item Downstream consolidation increases bargaining power, pressuring upstream firms to consolidate in response.
\item Technology shocks propagate through supplier relationships.
\item Efficiency gains in one stage of production create demand for reorganization at other stages.
\end{itemize}

\subsection*{1.6 Implications}
\begin{itemize}
\item Merger waves have a network structure; isolated industry studies miss this.
\item Merger policy (antitrust) should consider network effects of consolidation.
\item Industry networks are state variables for M\&A forecasting.
\end{itemize}

\subsection*{1.7 Caveats}
\begin{alertbox}
\begin{itemize}
\item I/O tables updated infrequently; may miss recent network changes.
\item Common shocks (macro) can drive simultaneous waves; identification requires lags.
\item Reverse causality: forward-looking consolidation in one industry may reflect anticipated consolidation in another.
\end{itemize}
\end{alertbox}

\section*{Round 2: Level 1 Fill-in}
\begin{enumerate}
\item Data: BEA \blank{1cm}-\blank{1cm} tables.
\item Sample: SDC mergers \blank{1cm}--\blank{1cm}.
\item Lagged \blank{2cm}/supplier waves predict focal-industry waves.
\item Spillover magnitude: \blank{1cm}--\blank{1cm}\% of SD.
\item Direction: \blank{2cm} more strongly propagates downstream.
\end{enumerate}

\section*{Round 3: Level 2 Prompts}
\begin{enumerate}
\item Describe the input-output data and link construction.
\item Report the main spillover findings.
\item Discuss mechanisms: bargaining, technology, efficiency.
\item Compare AH to Harford (2005).
\item What are the antitrust implications?
\end{enumerate}

\section*{Round 4: Minimal Scaffolding}
From a blank page: (i) question, (ii) I/O data, (iii) spillovers, (iv) mechanism, (v) policy.

\section*{Round 5: Blank Page}
Essay: merger-wave propagation through production networks.

\vspace{6pt}
\begin{drillbox}
\textbf{Speed Drill.} (1) Network data. (2) Sample. (3) Spillover size. (4) Upstream vs.\ downstream. (5) Mechanism.
\end{drillbox}

\section*{Quick Reference \& Answer Key}
\begin{answerbox}
\begin{itemize}
\item \textbf{Network.} BEA I/O tables.
\item \textbf{Finding.} Waves propagate via supplier-customer links.
\item \textbf{Size.} $\sim$20--30\% SD of focal activity.
\item \textbf{Lesson.} M\&A is networked; antitrust should notice.
\end{itemize}
\end{answerbox}

\begin{keybox}
\textbf{30-second exam pitch.} Ahern \& Harford (2014) test whether merger waves propagate across industries through supply-chain linkages. Using SDC mergers 1986--2010 and BEA input-output tables to measure industry customer/supplier ties, they estimate industry-year panels in which lagged merger activity in customer and supplier industries predicts focal-industry waves. A one-SD increase in customer/supplier merger activity raises focal merger activity by $\sim$20--30\%, with upstream-to-downstream propagation stronger than the reverse and effects decaying after 1--3 years. Mechanisms include consolidation-induced bargaining-power shifts and technology propagation. The paper establishes a network view of M\&A waves and complements Harford (2005) --- within-industry waves are real but partially induced by connected-industry consolidation --- with antitrust and industry-forecasting implications: merger policy must recognize that mergers in one industry set off consolidation in connected ones.
\end{keybox}

\end{document}
